% STATUS: VERIFIED / AUDITED
% FINAL DISPOSITION: VERIFIED / AUDITED AFTER MINOR REVISION
% AUDIT FILE: audits/A08_matched_bilinear_analysis.tex
% ORIGINAL AUDIT VERDICT:
%   MB FOUNDATION VERIFIED
%   MB SPECTRAL VERIFIED
%   SAME-TUPLE COMPARISON MINOR REVISION
%   PARITY/OBSTRUCTION VERIFIED
%   OVERALL MINOR REVISION
% AUDIT CLOSURE: A08 E01 applied verbatim; no other body change.
% ORIGINAL AUDITED VERSION: P08 v1.0; 2026-10-03; 39287 bytes; 846 lines
% ORIGINAL AUDITED FULL-FILE SHA-256: d21a77de84df4c40c2995d30ab740130795adca6c72f52fef9453ad5ddfb18ac
% ORIGINAL AUDITED BODY SHA-256: 3c4393e8a8285f946a67634addbd71d617985f7f153f4963fdcbde90ca123a0d
% CURRENT VERSION: P08 v1.1; 2026-10-03
% CURRENT BODY SHA-256: 67909444bd10dde10f2784db560e7ffb4165f34b63c0a72a356a35d38ba4cc85
% CURRENT HASH: full-file SHA-256 recorded in project_manifest.md and
%   audit_closure_report.txt, avoiding a self-referential file hash.
% BODY HASH SCOPE: bytes from first line-anchored section command through EOF.
% CLOSURE DATE: 2026-10-03
% SOLE ISSUE E01: original lines 534-535; n=8 MB representative 2, labels 2,6.
% PRESERVATION: all theorem statements, proofs, formulas, witnesses unchanged.
% STATUS NOTE: historical DRAFT/audit-pending prose in the body is retained;
%   this header and current manifest govern lifecycle status.
% VALIDATION PROVENANCE: A08 located and exactly reproduced the original
%   P08 validation companion; no validation checker rerun during closure.
% DOWNSTREAM: P07 cleared to begin, but not begun; later modules untouched.
% DEPENDS ON: P01-P05 VERIFIED / AUDITED
% MODELS: canonical DM and primary matched MB
% PRIMARY COMPARISON: same positive coefficient tuple, own positive synchronous
%                     equilibrium in each model
% NOVELTY ROLE: MAIN THEOREM TARGET if substantive qualitative contrast proved
% NUMERICAL DEPENDENCIES: none
% NEXT DEPENDENCIES: P07, P09/P10, P11
% BASELINE: cyclic_network_project_P05_audited.zip; 98 authenticated members.
% EXCLUSIONS: no later module, first-crossing, Hopf, l1 or six-term theorem.
% CLAIMS: C053-C063 as scoped in registry; P08-D1-P08-D9. C064/C065 deferred.

\section{P08: Matched bilinear analysis and exact comparison}
\label{sec:P08}
\noindent\textbf{Status: DRAFT; independent A08 audit required.}
All coefficients are strictly positive, $n\ge5$, $1\le q<n$, and indices
are modulo $n$. The comparator is exactly the five-parameter model
\eqref{eq:canonical-mb}. The same tuple $(a,k,c,g,h)$, $n,q$ and time
variable are used in DM and MB; each equilibrium amplitude is derived
for its own equation. Only current audited P01--P05 are imported.
Their historical DRAFT body sentences are superseded by their current
headers and audit-closure records.

The substantive results are open same-tuple stability contrasts in both
directions at $n=7$, a universal forward contrast, and realization of
every ordered pair of reachable dominant representatives with both
models stable. None is a first-instability, Hopf, criticality, or
literature-priority theorem. Here ``strictly stable'' means Hurwitz;
``unstable'' means a strictly positive spectral abscissa. Zero abscissa
is treated separately. Selection maximizes over all labels, including
zero, retaining every maximal tie.

\subsection{MB local flow, synchrony, and synchronous growth}
\label{slot:P08-1}
Put $S=g+h>0$ and $M=k+S-c$; the subspace remains denoted by $\Sync$.
\begin{theorem}[MB foundation]
\label{thm:P08-foundation}
For arbitrary canonical $(n,q)$ the following hold.
\begin{enumerate}[label=(\roman*),leftmargin=*]
\item MB has a unique maximal classical solution through every point of
$\R^n$. Every coordinate face is invariant. Nonnegative initial data
stay nonnegative, and initially positive coordinates remain strictly
positive on the maximal interval. No uniform asymptotic lower bound
is meant by this finite-time positivity statement.
\item $\Sync$ is invariant, with exact restriction
\begin{equation}
 \dot s=as-Ms^2.\label{eq:P08-scalar}
\end{equation}
The real synchronous equilibria are $0,a/M$ if $M\ne0$, and only $0$
if $M=0$. A positive synchronous equilibrium exists if and only if
$M>0$, and is uniquely $y^*=a/M$. For $M<0$ the nonzero root is negative.
\item For $s_0>0,M>0$, synchronous dynamics are global,
$0<s(t)\le\max(s_0,a/M)$ and $s(t)\to a/M$. For $M=0$ they are
$s(t)=s_0e^{at}$, global and unbounded. For $M<0$ they blow up at
\begin{equation}
 t_b=\frac1a\log\left(1+\frac{a}{|M|s_0}\right).
 \label{eq:P08-blowup}
\end{equation}
These classify nonnegative synchronous dynamics, not all MB trajectories.
\end{enumerate}
\end{theorem}
\begin{proof}
The polynomial field is smooth and locally Lipschitz on $\R^n$, so
local existence, uniqueness, continuous dependence and maximal
continuation apply. On the coordinate face $y_i=0$, direct substitution
in the original five terms gives $F_i^{\MB}(y)=0$, even when other
coordinates are negative. Along any solution the integrating-factor
identity is
\begin{equation}
 y_i(t)=y_i(0)\exp\int_0^t
 [a-ky_i(r)+cy_{i-1}(r)-gy_{i+q}(r)-hy_{i-q}(r)]\,dr.
 \label{eq:P08-coordinate}
\end{equation}
The integrand is continuous on each compact subinterval of the maximal
interval. The identity holds also for zero initial coordinates and
proves sign preservation and face invariance without assuming global
existence.

At $y=s\one$, each component of the original field is
$as-ks^2+cs^2-gs^2-hs^2=as-Ms^2$. The scalar solution lifts to a
vector solution, uniqueness identifies the two, and continuation makes
their maximal intervals coincide. Zeros of $s(a-Ms)$ give the complete
synchronous equilibrium list, retaining zero before division. For
$s_0>0$, solve $(1/s)'=-a/s+M$ to obtain
\begin{equation}
 s(t)=\frac{s_0e^{at}}{1+(Ms_0/a)(e^{at}-1)}.\label{eq:P08-logistic}
\end{equation}
For $M>0$ the denominator stays positive and the bound and limit follow;
for $M=0$ it is one; for $M<0$ its first zero is exactly
\eqref{eq:P08-blowup}, approached from above with positive numerator.
The zero initial state stays zero in every case.
\end{proof}

\subsection{Sufficient full-system bounds and their scope}
\label{slot:P08-3}
\begin{proposition}[Maximum comparison]
\label{prop:P08-bounds}
For nonnegative MB solutions and $Y=\max_i y_i$,
\begin{equation}
 D^+Y\le aY-(k-c)Y^2.\label{eq:P08-maximum}
\end{equation}
Thus $k\ge c$ suffices for global nonnegative existence. If $k>c$,
put $R=a/(k-c)$. Then $Y(t)\le\max\{Y(0),R\}$, $[0,R]^n$ is
invariant and uniformly attracts bounded subsets, and every
$[0,R+\varepsilon]^n$, $\varepsilon>0$, absorbs bounded subsets in
finite time. MB is therefore dissipative in this sufficient regime.
The comparison alone asserts neither finite-time absorption into the
exact cube $[0,R]^n$ nor boundedness when $k=c$.
\end{proposition}
\begin{proof}
Apply the generic finite-maximum identity of Lemma~\ref{lem:P01-max}.
At any maximizing coordinate $y_i=Y$, all neighbors lie in $[0,Y]$,
so $cy_i y_{i-1}\le cY^2$ and both inhibitory terms are nonpositive.
This includes tied maxima and coincident offsets. For $k\ge c$, the
inequality gives $Y(t)\le Y(0)e^{at}$, excluding finite-time escape;
polynomial continuation yields global existence. For $k>c$, logistic
comparison bounds $Y$ and proves invariance. A logistic solution
starting from the supremum of any bounded set of initial maxima
converges to $R$ and enters $[0,R+\varepsilon]$ in finite time,
proving the uniform attraction and absorption assertions.
\end{proof}
Consequently $M>0$ is necessary for boundedness of every nonnegative
MB trajectory, and $M\ge0$ is necessary for global existence of every
such trajectory, by the synchronous counterexamples. No converse is
inferred. In particular existence of positive synchrony is not used to
assert full-system boundedness. The sufficient condition $k>c$ is not
necessary, as an explicit example below will prove.

\subsection{Homogeneous MB Jacobian and arbitrary-offset spectrum}
\label{slot:P08-4}
At $y=s\one$, differentiate the original field before using equilibrium.
None of $i-1,i+q,i-q$ equals $i$ in the canonical range, so the diagonal is
\[
 d_{\MB}(s)=a-2ks+cs-gs-hs=a-(k+M)s.
\]
The three off-diagonal contributions are $cs,-gs,-hs$. Therefore
\begin{equation}
 DF^{\MB}(s\one)=[a-(k+M)s]I+s(cP-gP^{-q}-hP^q).
 \label{eq:P08-state-J}
\end{equation}
In full, with row $i$ and column $j$, the general derivative entry is
\[
 [a-2ky_i+cy_{i-1}-gy_{i+q}-hy_{i-q}]\delta^{(n)}_{ji}
 +cy_i\delta^{(n)}_{j,i-1}-gy_i\delta^{(n)}_{j,i+q}
 -hy_i\delta^{(n)}_{j,i-q}.
\]
Coincident columns add; none is overwritten. At $y^*=a/M>0$ use
$a=My^*$ only now, giving $d_{\MB}=-ky^*$.

\begin{theorem}[Full MB spectrum, normality, and uniform direction]
\label{thm:P08-spectrum}
For $M>0$, positive MB synchrony has
\begin{align}
 J_{\MB}&=y^*[-kI+cP-gP^{-q}-hP^q],\label{eq:P08-J}\\
 \lambda_m^{\MB}&=y^*[-k+ce^{-\ii\theta_m}-ge^{\ii q\theta_m}
                           -he^{-\ii q\theta_m}],\label{eq:P08-spectrum}\\
 \Ree\lambda_m^{\MB}&=y^*[-k+c\cos\theta_m-S\cos(q\theta_m)],
 \label{eq:P08-real}\\
 \Imm\lambda_m^{\MB}&=y^*[-c\sin\theta_m+(h-g)\sin(q\theta_m)].
 \label{eq:P08-imag}
\end{align}
The $n$ Fourier labels give the full spectrum with algebraic multiplicity;
$J_{\MB}$ is normal and unitarily diagonalized by the project's DFT.
In particular
\begin{equation}
 \boxed{\lambda_0^{\MB}=-a<0.}\label{eq:P08-uniform}
\end{equation}
Any linear instability at positive MB synchrony involves a nonuniform
mode. This does not classify a bifurcation as Hopf or steady.
\end{theorem}
\begin{proof}
The preceding differentiation proves the matrix identity. In the verified
convention $P^\ell v^{(m)}=e^{-\ii\ell\theta_m}v^{(m)}$. Applying
these powers yields the stated spectrum, including the negative
coefficient on the positive-phase $g$ term. Euler's formula gives the
real and imaginary parts; the imaginary contributions are
$-c\sin\theta_m,-g\sin(q\theta_m),+h\sin(q\theta_m)$.
The DFT is unitary by the generic geometric-sum proof in P04, whence
$J_{\MB}=\Fourier\diag(\lambda_m^{\MB})\Fourier^*$, proving the
full multiplicity statement and normality. Independently, $J_{\MB}$
and its adjoint are sums of mutually commuting powers of the unitary
shift, so they commute.

For the first check of the uniform value, Fourier substitution at $m=0$
gives $y^*(-k+c-g-h)=-My^*=-a$. For a second independent check,
the derivative of the scalar restriction is $a-2Ms$, equal to $-a$
at $s=a/M$. Its tangent vector is $\one$, so it is precisely the
uniform eigenvalue. Its strict negativity proves the instability claim.
\end{proof}

The complete offset-collision checks, using P03/P04 only for their
index geometry, are
\[
\begin{array}{c|l}
q=1&J_{\MB}/y^*=-kI+(c-h)P-gP^{-1}\\
q=n-1&J_{\MB}/y^*=-kI+(c-g)P-hP^{-1}\\
q=n/2\ (n\text{ even})&J_{\MB}/y^*=-kI+cP-SP^{n/2}.
\end{array}
\]
Their imaginary parts are respectively
$y^*(-c+h-g)\sin\theta_m$, $y^*(-c+g-h)\sin\theta_m$, and
$-y^*c\sin\theta_m$. The last real restraint term is $-y^*S(-1)^m$.
Each also follows from the master spectrum, with cancellation retained.
For even $n$, the singleton is
$\lambda_{n/2}^{\MB}=y^*[-k-c-S(-1)^q]$; the uniform value is
unchanged in every collision case. Label conjugacy and semisimplicity
follow from the real circulant matrix and its unitary diagonalization;
a pair may have zero imaginary part or share an eigenvalue with other
labels. No numerical-eigenvalue simplicity is inferred from label counts.

\begin{corollary}[Exact stability and another sufficient global bound]
\label{cor:P08-stability}
Define
\begin{equation}
 K_{n,q}(c,S)=\max_{0\le m<n}[c\cos\theta_m-S\cos(q\theta_m)].
 \label{eq:P08-K}
\end{equation}
At positive MB synchrony, strict stability is equivalent to $k>K_{n,q}$,
and instability to $k<K_{n,q}$. At equality the normal linear system
is Lyapunov stable but not asymptotically stable; nonlinear stability is
undecided here. Hyperbolicity requires every real part to be nonzero,
not merely a nonzero maximum.

Moreover $k>K_{n,q}$ itself implies $M>0$, global boundedness and
dissipativity for nonnegative MB data. This is another sufficient bound,
not a necessary condition or an identification with $k>c$.
\end{corollary}
\begin{proof}
The abscissa is $y^*(K_{n,q}-k)$ and unitary diagonalization gives
$\|e^{tJ_{\MB}}\|_2=e^{t\sa(J_{\MB})}$, proving the linear claims.
For the bound put $\delta=k-K_{n,q}>0$, $V=\sum_i y_i$. Summing the
original equations yields
\begin{equation}
 V'=aV+y^{\mathsf T}[-kI+cP-gP^{-q}-hP^q]y
 \le aV-\delta\|y\|_2^2\le aV-\frac\delta nV^2.
 \label{eq:P08-sum}
\end{equation}
The symmetric part of the displayed matrix has eigenvalues
$-k+c\cos\theta_m-S\cos(q\theta_m)\le-\delta$, and its skew part
contributes zero to the real quadratic form. Logistic comparison bounds
$V$ before any maximal endpoint and proves global continuation.
The simplex $V\le na/\delta$ attracts bounded subsets, and every
strictly larger one absorbs them. Finally $K_{n,q}\ge c-S$ by its
zero label, so $M>0$.
\end{proof}
For example, for $q=2$ the tuple $(9/2,7,8,1,1)$ has $K_{n,2}=6$
for every $n\ge5$: $2+8u-4u^2$ uniquely maximizes at $u=1$ on the
grid. MB is dissipative although $k=7<c=8$. Its own root is $y^*=9/2$,
whereas DM's unique positive root is $x^*=1$. Direct original residuals
are $9/2-8+7/2=0$ for DM after division by $x^*$, and
$9/2-(9/2)=0$ for MB after division by $y^*$. This refutes necessity
of $k>c$. Conversely, $k>c$ need not imply spectral stability when
$S$ is large, so the sufficient bounds have different scopes.

\subsection{Exact step-two MB selected-mode rule}
\label{slot:P08-2}
Set $q=2$ and $u_m=\cos\theta_m$. Since $\cos(2\theta)=2u^2-1$,
\begin{align}
 p(u_m):=\Ree\lambda_m^{\MB}/y^*&=-k+S+cu_m-2Su_m^2,
 \label{eq:P08-quadratic}\\
 p(u)&=-k+S+\frac{c^2}{8S}-2S(u-v_{\MB})^2,\nonumber\\
 v_{\MB}&=\frac{c}{4S}=\frac1{\chi_{\MB}},\qquad
 \chi_{\MB}=\frac{4S}{c}>0.\label{eq:P08-vertex}
\end{align}
This is always strictly concave, $p''=-4S<0$. The unconstrained vertex
is $v_{\MB}$; on $[-1,1]$ the continuous maximizer is
$\min(1,v_{\MB})$. Its departure from $1$ need not switch a discrete mode.

Use P05's notation
\[
 N=\lfloor n/2\rfloor,\quad u_j=\cos(2\pi j/n),\quad
 \beta_j=(u_j+u_{j+1})/2,\quad J=\lfloor(n+1)/4\rfloor,
 \quad\alpha_j=1/\beta_j\ (0\le j<J).
\]
P05 proves $\beta_j=\cos(\pi/n)\cos((2j+1)\pi/n)$ and exactly
$\beta_0,\ldots,\beta_{J-1}>0$. Its nearest-grid-cell argument,
midpoint ordering, parity and label bookkeeping are generic finite-grid
facts. Its DM realization theorem is not an MB realization proof.

\begin{theorem}[MB cells, all ties, reachability, and cell stability]
\label{thm:P08-selection}
On $M>0$, the complete MB representative selection is
\begin{equation}
\begin{array}{c|c}
0<\chi_{\MB}<\alpha_0&\{0\}\\
\alpha_{j-1}<\chi_{\MB}<\alpha_j\ (1\le j<J)&\{j\}\\
\chi_{\MB}>\alpha_{J-1}&\{J\}\\
\chi_{\MB}=\alpha_j\ (0\le j<J)&\{j,j+1\}.
\end{array}\label{eq:P08-cells}
\end{equation}
The middle row is empty for $J=1$. Equivalently, closed vertex cells are
$[\beta_0,\infty)$, $[\beta_j,\beta_{j-1}]$ ($1\le j<N$), and
$(-\infty,\beta_{N-1}]$, intersected with $v_{\MB}>0$.
Every listed strict cell and tie is realizable with $M>0$; the exact
reachable representatives are $0,\ldots,J$.

A nonzero selected representative gives $\{j,n-j\}$, while zero is a
singleton. A tie involving zero gives three labels, other reachable
ties four. No nonadjacent or triple representative maximum occurs.
The even Nyquist singleton is excluded. These are label counts, not
numerical-eigenvalue multiplicities or frequency assertions.
On the cell for $j$ strict stability is exactly
\begin{equation}
 \boxed{k>S+cu_j-2Su_j^2.}\label{eq:P08-cell-stability}
\end{equation}
Either selected representative gives the same bound at a tie.
Reverse strict inequality gives instability and equality gives zero
abscissa. In the uniform cell, including a maximal tie with zero,
$M>0$ already makes the maximum $-a<0$.
\end{theorem}
\begin{proof}
Square completion reduces selection to minimizing $|u_j-v_{\MB}|$.
Apply the generic parts of Theorems~\ref{thm:P05-selection} and
\ref{thm:P05-reachable}. An independent sign check is
\[
 p(u_j)-p(u_{j+1})=c(u_j-u_{j+1})(1-\chi_{\MB}\beta_j),
\]
confirming the switch orientation. For MB-specific realizability,
prescribe any $v>0$, take $g=h=1/2$, $c=4v$,
$k>\max(0,c-1)$ and any $a>0$. Then $M>0$ and $v_{\MB}=v$.
Each required strict cell and positive midpoint can be attained; strict
cells are open by continuity. The grid theorem gives the labels,
exclusions and ties. Taking the maximum of the derived quadratic proves
the exact cell-stability test.
\end{proof}
The precise parity effects follow from the same generic grid calculation:
\[
\begin{array}{c|c|c|c}
n&J&u_J&\text{nearest representatives at formal vertex }0\\\hline
4\ell&\ell&0&\ell\\
4\ell+1&\ell&\sin(\pi/(2n))&\ell\\
4\ell+2&\ell&\sin(\pi/n)&\ell,\ell+1\\
4\ell+3&\ell+1&-\sin(\pi/(2n))&\ell+1.
\end{array}
\]
Every finite positive tuple has $v_{\MB}>0$; the formal tie for
$n=4\ell+2$ is unattainable. Its two real parts differ by
$2y^*c\sin(\pi/n)>0$. Only the gap divided by $y^*$ is independent
of $S$; the unnormalized gap need not stay away from zero in parameter
limits. At fixed $n$, sufficiently large finite $\chi_{\MB}$ uniquely
selects $J$ in every parity class. No large-$n$ theorem is asserted.

\subsection{Common domain, named branches, and structural comparison}
\label{slot:P08-5}
For DM set $A=k+h$, $\mu=a+c-g$, $b=a-A$, and $D=b^2+4A\mu$.
By audited Theorem~\ref{thm:P02-positive}, the full common existence set is
\begin{equation}
\begin{split}
 \mathcal C=\{(a,k,c,g,h)>0:\ &M>0,\quad
 [\mu>0\ \text{or}\ (\mu=0,b>0)\\
 &\text{or}\ (\mu<0,b>0,D\ge0)]\}.
\end{split}\label{eq:P08-common}
\end{equation}
It is not entirely an open set. The branch-labelled comparison pairs
a tuple with one of its allowed positive DM roots
$x_\pm=(b\pm\sqrt D)/(2A)$ and the sole MB root $a/M$.
Most constructions use the open domain
\begin{equation}
 \mathcal C_+=\{(a,k,c,g,h)>0:M>0,\ \mu>0\},\qquad x^*=x_+.
 \label{eq:P08-common-simple}
\end{equation}
DM then has exactly one positive synchronous equilibrium; its scalar
uniform derivative is $-x^*\sqrt D/(1+x^*)<0$ by P04. A separate
construction below explicitly follows $x_-$ on a two-positive-root domain.

At a specified positive DM equilibrium put $s=1+x^*$. Quote P03--P05:
\begin{equation}
\begin{gathered}
 d=-kx^*-c(1+2x^*)/s^2,\quad C=c/s,\quad G=gx^*/s^2,\quad H=hx^*,\\
 r_m^{\DM}=d+Cu_m+(G-H)(2u_m^2-1),\\
 \chi_{\DM}=\frac{4(H-G)}C
 =\frac{4x^*[h(1+x^*)^2-g]}{c(1+x^*)},\qquad
 \chi_{\MB}=\frac{4(g+h)}c.
\end{gathered}\label{eq:P08-DM-input}
\end{equation}
DM's parameter can be negative, zero or positive; its positive-vertex
interpretation applies only if $h>g/s^2$. P05's full selection rule uses
$\chi_{\DM}<\alpha_0$ for every uniform case. At fixed $k,c,g,h$,
MB selection is independent of $a$, whereas DM selection can vary with
$a$ through the equilibrium branch.

\begin{proposition}[Structural functional-form distinctions]
\label{prop:P08-structural}
MB real dispersion for $q=2$ is always strictly concave in $u$.
DM dispersion is strictly convex, linear, or strictly concave according
as $h<g/(1+x^*)^2$, $h=g/(1+x^*)^2$, or $h>g/(1+x^*)^2$.
When $h\le g/(1+x^*)^2$, DM uniquely selects zero. All three
curvature regimes occur in $\mathcal C_+$.

MB has invariant coordinate faces. DM has
$F_i^{\DM}|_{x_i=0}=cx_{i-1}$ and propagates positivity from any nonzero
nonnegative initial state by P01. MB's uniform eigenvalue at positive synchrony is
always $-a$. DM's is $-x^*\sqrt D/s$ on $x_+$,
$+x^*\sqrt D/s$ on a distinct positive $x_-$, and zero at a positive
double root.

DM is globally nonnegative and dissipative for the entire canonical
positive class by P01. MB has the sufficient dissipative regimes above,
but global unbounded synchronous growth at $M=0$ and finite-time
synchronous blow-up at $M<0$. Thus the DM functional form is
structurally dissipative over the entire canonical class, whereas MB is
not. The growth/blow-up examples are outside $\mathcal C$.
\end{proposition}
\begin{proof}
The second derivatives are $-4y^*S<0$ and $4(G-H)$, respectively;
use P05 for uniform selection. To realize the three curvature signs,
set $x^*=1$, $c=g=4$, $k=1$, $a=1+h$, with $h=1/2,1,2$.
Then $\mu=a>0$, $M=1+h>0$, $y^*=1$, $G-H=1-h$, and both
original scalar residuals vanish. The face and global statements quote
P01 and the MB proofs above. The uniform DM values result from inserting
P02's roots in \eqref{eq:P04-equilibrium-zero-mode}. No full-system
boundedness is inferred from MB equilibrium existence.
\end{proof}
\begin{remark}[Meaning and limits of matching]
Matching fixes $n,q$, the five positive coefficients, directed offsets
and coefficient roles, and the original time variable. It does not
calibrate equilibrium amplitude, local derivatives, empirical behavior,
or nonlinear response. The successor derivative is positive for DM and
negative for MB at positive synchrony, despite the negative term in both
fields. No biological or quantitative-superiority claim follows.
For a fixed source $z>0$, as receiver $r\downarrow0$, the predecessor
inputs satisfy $cz/(1+r)\to cz$ for DM and $crz\to0$ for MB.
This last assertion is only a limit under the explicitly fixed-source
convention.
\end{remark}

\subsection{Every ordered dominant-mode pair with both models stable}
\begin{theorem}[Independent realization of the two selection parameters]
\label{thm:P08-pair-realization}
For each fixed $n\ge5,q=2$, every ordered pair
$(j_{\DM},j_{\MB})\in\{0,\ldots,J\}^2$ occurs with unique selected
representatives and both equilibria strictly stable on a nonempty open
subset of $\mathcal C_+$. In particular reverse uniform/nonuniform
selection is possible; distinct nonuniform/nonuniform pairs are possible
when $J\ge2$, beginning at $n=7$. This is not a full stability partition.
\end{theorem}
\begin{proof}
Choose positive parameters in the desired strict cells; explicit choices
for representative $j$ are
\[
 z_0=\alpha_0/2,\quad
 z_j=(\alpha_{j-1}+\alpha_j)/2\ (1\le j<J),\quad
 z_J=2\alpha_{J-1}.
\]
When $J=1$ use only the first and last expressions. Put
$\eta=z_{j_{\DM}}$, $\zeta=z_{j_{\MB}}$, and define
\begin{equation}
\begin{gathered}
 c=4,\quad S=\zeta,\quad x=1+\eta/\zeta,\quad s=1+x,\\
 h=\frac{\zeta+s\eta/x}{s^2+1},\quad g=\zeta-h,\\
 k>\max\{c+S,\ (c-g)/(xs)-h\},\qquad
 a=(k+h)x-(c-g)/s.
\end{gathered}\label{eq:P08-pair-construction}
\end{equation}
Here $xs>\eta/\zeta$, so $0<h<S$ and $g,h>0$. The bound on $k$
gives $a>0$ and $M>0$. Further,
$\mu=x[k+h+(c-g)/s]>0$, since $k>S>g$. Substitution into the
\emph{original} vector fields gives
\begin{equation}
 F_i^{\DM}(x\one)=x[a-(k+h)x+(c-g)/s]=0,\qquad
 F_i^{\MB}(y^*\one)=y^*(a-My^*)=0.
 \label{eq:P08-residual-general}
\end{equation}
Algebra gives $\chi_{\MB}=\zeta$ and
$x(hs^2-g)/s=\eta$, hence $\chi_{\DM}=\eta$.
For all MB labels $\Ree\lambda_m^{\MB}/y^*\le-k+c+S<0$.
For DM the audited formula gives
\begin{equation}
 \frac{r_m^{\DM}}x
 \le-k-\frac{c}{s^2}+\frac{g}{s^2}+h<-k+S<0.
 \label{eq:P08-DM-bound}
\end{equation}
Indeed the two $c$ contributions after division by $x$ are at most
$-c/s^2$, and $|g/s^2-h|\le g/s^2+h$. This proves all claimed
strict inequalities. Openness in all physical parameters, not just this
displayed family, is proved below.
\end{proof}
A rational reverse-selection witness at $n=8$ is
\begin{equation}
 (a,k,c,g,h)=(113/4,9,8,1,1),\quad x^*=3,\quad y^*=113/12.
 \label{eq:P08-reverse-selection-example}
\end{equation}
Here $M=3$, $\mu>0$, $\chi_{\MB}=1$ selects $0$, while
$\chi_{\DM}=45/8>2\sqrt2=\alpha_1$ selects $2$.
The abscissae are $-443/16$ for DM and $-113/4$ for MB.
Original residuals after division by the respective amplitudes are
$113/4-10\cdot3+7/4=0$ and $113/4-3(113/12)=0$.

\subsection{DM stable and MB unstable for every size}
\begin{theorem}[Universal same-tuple opposite stability]
\label{thm:P08-forward}
For each fixed $n\ge5,q=2$, a nonempty open subset of $\mathcal C_+$
has DM strictly stable with unique dominant representative $0$, while
MB is unstable with a unique nonzero dominant representative. A single
rational tuple works for every dimension:
\begin{equation}
 \boxed{(a,k,c,g,h)=(1/4,1/8,1,1,1/8).}
 \label{eq:P08-forward-witness}
\end{equation}
At this tuple the independently computed amplitudes happen to be
$x^*=y^*=1$; amplitude equality is not a comparison assumption.
\end{theorem}
\begin{proof}
All coefficients are positive, $M=1/4$, $\mu=1/4$. At amplitude one
the original DM residual is $1/4-1/8+1/2-1/2-1/8=0$ and the MB
residual is $1/4-1/8+1-1-1/8=0$. P02 ensures the unique positive DM
root. Its coefficients are $d=-7/8,C=1/2,G=1/4,H=1/8$.
Thus $G-H=1/8>0$, zero is strictly dominant and
$r_0^{\DM}=-1/4$; every DM real part is negative.
For MB,
\[
 S=9/8,\qquad p(u)=1+u-\frac94u^2,\qquad
 v_{\MB}=2/9,\qquad\chi_{\MB}=9/2.
\]
The parity formulas give $-1/4<u_J\le1/2$.
In the only negative case $n\equiv3\pmod4$, $n\ge7$ and
$|u_J|=\sin(\pi/(2n))<\pi/14<1/4$. The other cases give the
upper bound, attained at $n=6$. Hence the actual finite label $J$ satisfies
\[
 p(u_J)\ge1-\frac14-\frac94\left(\frac12\right)^2
 =\frac3{16}>0.
\]
MB is unstable, and its negative uniform eigenvalue excludes zero from
the selected set. There is no representative tie: a tie would force
$\beta_j=2/9$, but
$4\beta_j=2\cos(2\pi j/n)+2\cos(2\pi(j+1)/n)$ is an algebraic
integer, being a sum of roots of unity and their inverses. The rational
number $8/9$ is not an integer, whereas every rational algebraic integer
is an integer (apply the rational-root argument to a monic integer
polynomial). This excludes a tie. Strict gaps and branch regularity
give openness by the argument below.
\end{proof}
For $n=5,6,7$ this witness selects MB representative $1$; for $n=8$
it selects representative $2$. Other sizes follow the exact finite-grid
rule and need not retain either label.
The instability proof uses a finite Fourier label rather than only a
positive continuous maximum.

\subsection{Reverse stability: parity obstruction and constructive regions}
For a positive DM equilibrium write $x=x^*$, $s=1+x$,
$u=\cos\theta_m$, $w=-\cos(2\theta_m)$. The audited DM formula and
the derived MB formula give
\begin{equation}
 \frac{r_m^{\DM}}x=-k+(h-g/s^2)w
       -\frac{c(1+2x-su)}{xs^2},\qquad
 \frac{r_m^{\MB}}{y^*}=-k+Sw+cu.
 \label{eq:P08-modal-comparison}
\end{equation}
Notice $1+2x-su\ge x>0$ for $u\le1$.
\begin{theorem}[Reverse stability on the unique positive branch]
\label{thm:P08-reverse}
On $\mathcal C_+$ with $q=2$, a tuple with DM unstable and MB strictly
stable exists if and only if $n\equiv3\pmod4$. For every such fixed
size there is a nonempty open region realizing this contrast, with both
models uniquely selecting $J$. For all other sizes, MB strict stability
implies DM strict stability on $\mathcal C_+$.
\end{theorem}
\begin{proof}[Proof of the obstruction]
On $\mathcal C_+$ the uniform DM eigenvalue is strictly negative.
If $G-H\ge0$, P05 makes it largest, so DM is strictly stable.
Otherwise $h-g/s^2>0$. A nonnegative DM real part in
\eqref{eq:P08-modal-comparison} forces $w>0$. If its cosine $u\ge0$,
\begin{equation}
 \frac{r_m^{\MB}}{y^*}-\frac{r_m^{\DM}}x
 =g(1+s^{-2})w+cu+\frac{c(1+2x-su)}{xs^2}>0.
 \label{eq:P08-modal-gap}
\end{equation}
That MB mode is therefore positive. Unless $n\equiv3\pmod4$, all
reachable DM representatives have nonnegative cosines by P05. Any
nonnegative DM maximum would contradict MB strict stability. This also
excludes a zero DM abscissa and proves the strict implication.
\end{proof}
\begin{proof}[Proof of the remaining-size realization]
Let $n\equiv3\pmod4$, $\tau=\sin(\pi/(2n))$, so $0<\tau<1/4$,
$u_J=-\tau$, and $W=1-2\tau^2>7/8$.
Also $u_{J-1}=\sin(3\pi/(2n))=3\tau-4\tau^3$, giving
$\beta_{J-1}=\tau W$. Set
\begin{equation}
\begin{gathered}
 h=1,\quad c=\tau,\quad g=\tau^2/4,\quad
 k=1-5\tau^2/2,\quad x=2/\tau,\quad s=1+x,\\
 a=(k+1)x-(\tau-\tau^2/4)/s.
\end{gathered}\label{eq:P08-reverse-family}
\end{equation}
Every coefficient is positive: $k>27/32$, $x>8$, the subtracted
term in $a$ is less than $1/4$, and $(k+1)x>8$.
Further $M=2-\tau-9\tau^2/4>0$, and $\mu>0$ since $c>g$.
The original residuals vanish by \eqref{eq:P08-residual-general}.

For MB, $v_{\MB}=\tau/(4+\tau^2)<\tau/4<\beta_{J-1}$ and
$\beta_J<0$, so it selects $J$ strictly. Its maximum divided by $y^*$ is
\[
 -k+(1+\tau^2/4)W-\tau^2=-\tau^2(1/4+\tau^2/2)<0.
\]
For DM the $J$-mode divided by $x$ is
\[
 \frac{r_J^{\DM}}x=\frac{\tau^2}{2}-\frac{gW}{s^2}
              -\frac{\tau(1+2x+s\tau)}{xs^2}.
\]
Since $1+2x+s\tau\le(2+\tau)s$, $s>x$, and $x=2/\tau$,
the two subtracted terms are less than
\[
 \frac{g+\tau(2+\tau)}{x^2}<\frac{37}{256}\tau^2.
\]
Here $g<1/64$ and $\tau(2+\tau)<9/16$. Thus
$r_J^{\DM}/x>91\tau^2/256>0$.
DM is concave, and its vertex obeys
\[
 v_{\DM}=\frac{\tau}{4x(s-g/s)}<\frac{\tau}{256}
      <\tau W=\beta_{J-1},
\]
since $x>8$ and $s-g/s>8$. So DM also strictly selects $J$.
These are strict canonical common-domain inequalities; openness is
established below.
\end{proof}

\paragraph{Rational reverse-stability witness.}
At $n=7$ a compact exact witness is
\begin{equation}
 (a,k,c,g,h)=(20301/1100,17/20,1/2,1/100,1),\quad
 x^*=10,\quad M=34/25,\quad y^*=20301/1496.
 \label{eq:P08-reverse-rational}
\end{equation}
Original residuals divided by their amplitudes are
$20301/1100-(37/20)10+(49/100)/11=0$ for DM and
$20301/1100-(34/25)(20301/1496)=0$ for MB.
All coefficients and $\mu$ are positive.
The ordered seven-node cosines $u_1,u_2,u_3$ are the three roots of
$8u^3+4u^2-4u-1$: use
$1+2(\cos\theta+\cos2\theta+\cos3\theta)=0$ at nonzero seventh roots
and the Chebyshev identities. Rational isolating intervals are
\begin{equation}
 31/50<u_1<63/100,\quad -23/100<u_2<-11/50,\quad
 -91/100<u_3<-9/10.\label{eq:P08-seven-intervals}
\end{equation}
The endpoint signs of the cubic are respectively $(-,+),(+,-),(-,+)$.
Three disjoint sign changes exhaust its degree; cosine ordering assigns
the labels. Thus no decimal trigonometric value is a proof premise.
MB representatives $1,2,3$ have normalized real parts
$-17/20+u_1/2-101u_2/100$,
$-17/20+u_2/2-101u_3/100$, and
$-17/20+u_3/2-101u_1/100$.
The interval upper bounds are respectively $-3027/10000$,
$-409/10000$, $-19262/10000$, all negative; mode zero is $-34/25$.
DM instead satisfies
\[
 \frac{r_2^{\DM}}{10}
 >-\frac{17}{20}-\frac{21}{2420}-\frac{23}{22000}
      +\frac{12099}{12100}\frac9{10}
 =\frac{9729}{242000}>\frac1{25}>0.
\]
Both select $2$ strictly: $v_{\MB}=25/202<39/200<\beta_1$,
$v_{\DM}=55/48396<39/200<\beta_1$, and $\beta_2<0$.

\paragraph{The full domain requires explicit branch labels.}
The preceding obstruction cannot be extended to all of $\mathcal C$
without qualification. A lower-branch construction gives the exact
existence answer on that larger branch-labelled domain.
\begin{proposition}[Reverse contrast when both DM branches are allowed]
\label{prop:P08-lower}
If $4\mid n$ and MB is strictly stable, every positive DM synchronous
equilibrium is strictly stable. If $4\nmid n$, a nonempty open common
region has MB strictly stable and the specified lower positive DM
branch unstable with unique dominant representative zero.
\end{proposition}
\begin{proof}
When $4\mid n$, MB's label $n/4$ has normalized real part $-k+S$.
Its strict stability forces $k>S$. Bound \eqref{eq:P08-DM-bound},
which did not require $\mu>0$, makes every DM real part negative.

Otherwise put $\rho_n=\max_m[-\cos(2\theta_m)]$. One has
$0<\rho_n<1$: equality to one requires $4m=n$ modulo $2n$, impossible
when $4\nmid n$, and a label nearest $n/4$ has negative double-angle
cosine. Define
\begin{equation}
\begin{gathered}
 \varepsilon=(1-\rho_n)/24,\quad k=1,\quad c=h=\varepsilon,\quad
 g=(1+\varepsilon)^3+2\varepsilon,\\
 a=(1+\varepsilon)(1+2\varepsilon)
                         +\frac{\varepsilon}{1+\varepsilon}.
\end{gathered}\label{eq:P08-lower-family}
\end{equation}
All coefficients are positive and $M=1+g>0$. With $A=1+\varepsilon$,
the two DM roots are exactly
\[
 x_-=\varepsilon,\qquad
 x_+=\varepsilon+\frac{\varepsilon}{(1+\varepsilon)^2}.
\]
Direct substitution of $x_-$ in $a-Ax+(c-g)/(1+x)$ gives zero.
The second is checked by the quadratic factorization
$A(x-x_-)(x-x_+)$ of P02's equilibrium polynomial, or by substitution.
In particular $\mu=-A\varepsilon^2-\varepsilon^2/(1+\varepsilon)<0$,
$b>0$, $D>0$, and both roots are positive. On the named lower branch,
\[
 \lambda_0^{\DM}=\frac{\varepsilon^2}{(1+\varepsilon)^2}>0,
 \qquad G-H=\varepsilon\left(1+
                       \frac{2\varepsilon}{(1+\varepsilon)^2}\right)>0.
\]
Thus zero is uniquely dominant and unstable. For MB,
$S=(1+\varepsilon)^3+3\varepsilon<1+10\varepsilon$, since
$0<\varepsilon<1$, and every normalized real part obeys
\[
 -1+\varepsilon\cos\theta_m-S\cos(2\theta_m)
 \le-1+\varepsilon+S\rho_n
 <-\frac{13}{24}(1-\rho_n)<0.
\]
The MB original residual vanishes at $y^*=a/M$. Distinct simple
positive DM roots and all strict inequalities persist locally on the
specified lower branch. At $n=6$, $\rho_6=1/2$, so
$\varepsilon=1/48$ gives an entirely rational witness.
\end{proof}

\subsection{Mode/stability combinations and physical-parameter openness}
All five requested types of combination have realizations, without a
claim that each occurs in every dimension. Both stable with identical
or different modes follows from Theorem~\ref{thm:P08-pair-realization};
DM stable / MB unstable is universal by Theorem~\ref{thm:P08-forward};
the reverse has the branch and size restrictions just proved.
For both unstable with different dominant modes take $n=8$ and
\begin{equation}
 (a,k,c,g,h)=(16351/1100,1/2,3/2,1/100,1),\quad
 x^*=10,\quad M=1/100,\quad y^*=16351/11.
 \label{eq:P08-both-unstable}
\end{equation}
The original residuals divided by their amplitudes are
$16351/1100-15+(149/100)/11=0$ and
$16351/1100-(1/100)(16351/11)=0$. Coefficients and $\mu$ are positive.
MB's vertex $75/202$ lies strictly between $\sqrt2/4$ and
$(2+\sqrt2)/4$, so MB selects $1$ with normalized maximum
$-1/2+3\sqrt2/4>0$. DM's vertex $55/16132<\sqrt2/4$ selects $2$,
whose normalized maximum is
\[
 -\frac12-\frac{63}{2420}+\frac{12099}{12100}
 =\frac{2867}{6050}>0.
\]
The small value of $M$ is strictly positive. MB's large amplitude is
its own computed equilibrium, not an amplitude imported from DM.

\begin{proposition}[Openness in all five physical coefficients]
\label{prop:P08-open}
Every strict contrast and strict mode pair constructed here persists
on a nonempty open set in $(a,k,c,g,h)\in(0,\infty)^5$ for its stated
fixed size, offset, and DM branch. No openness assertion rests merely
on an interval inside a displayed lower-dimensional family.
\end{proposition}
\begin{proof}
On $\mathcal C_+$, $D>0$ and the unique positive root $x_+$ is smooth
in all five coefficients. Equivalently, the scalar equilibrium function
$B(x)=a-(k+h)x+(c-g)/(1+x)$ has
$B_x=-\sqrt D/(1+x_+)<0$, so the implicit function theorem applies.
MB's branch $a/M$ is smooth for $M>0$. Coefficient positivity,
$M>0$, and $\mu>0$ are strict open conditions. At fixed $n$ there
are finitely many spectral inequalities; each is strict and continuous.
A unique selected representative has a strictly positive gap to every
other representative, while its conjugate labels remain tied.
A sufficiently small intersection of these neighborhoods preserves all
branch, sign, stability and selection assertions.

For the lower-branch family use instead $\mu<0,b>0,D>0$.
Both positive roots are simple and the specified $x_-$ is smooth, with
$B_x=+\sqrt D/(1+x_-)>0$. The same continuity argument preserves
that branch and its strict conclusions. No full-system IFT, later
module, uniform neighborhood radius in $n$, or global branch
identification is used.
\end{proof}

\subsection{Interpretation, deferred gates, and internal checks}
\label{slot:P08-6}
The strongest completed DRAFT results are a universal forward stability
contrast, both stability directions at $n=7$ on the unique-positive-root
domain, every ordered pair of reachable modes with both models stable,
and exact parity/branch limitations on reverse stability. Identical
abstract reachable labels do not imply identical behavior at the same
physical tuple. The unrestricted reverse-stability hope fails on
$\mathcal C_+$ unless $n\equiv3\pmod4$; even allowing both DM
branches cannot produce it when $4\mid n$.
The rejected phrase ``MB is generically unstable'' is not reinstated:
MB has explicit open stable regions and an exact stability criterion.

For structural growth examples, $(a,k,c,g,h)=(1,1,3,1,1)$ gives
$M=0$, and $(1,1,4,1,1)$ gives $M=-1$. At $s_0=1$ the latter blows
up at $\log2$. DM has its P01 bounded flow for both. Its own positive
root is $1$ or $(\sqrt{33}-1)/4$, respectively, each satisfying the
original residual $x[a-(k+h)x+(c-g)/(1+x)]=0$.
MB has no positive synchronous root in either case; its origin has zero
residual, and the stated scalar trajectories have the exact original
MB derivative. These are not common-equilibrium stability comparisons.

\label{slot:P08-extra-1}
The six-term model is explicitly deferred. Existing source-provenance
records are retained, but no archived six-term equation is imported.
Full MB Hopf, first Lyapunov coefficient, crossing transversality, and
criticality comparisons are deferred as well. The spectra and open
regions are inputs for later P11--P13, not proofs of their hypotheses.
A selected nonuniform pair, positive real part, or candidate imaginary
crossing is insufficient for Hopf. P07, P09, P10 and all later proof
bodies are untouched. The frozen package already calls for this
comparison, so no scope amendment or novelty-priority promotion is needed.

\paragraph{Internal analytic and computational checks.}
The proofs independently check the original scalar reduction, coordinate
faces, reciprocal scalar solution, maximum inequality, pre-equilibrium
diagonal, Fourier phases, both uniform-eigenvalue derivations, quadratic
completion, midpoint orientation, and parameter constructions.
Each displayed exact equilibrium has an original-field residual check;
the general constructions have symbolic residual identities.
The separate validation-support package supplies an independent
term-by-term field/Jacobian implementation, direct matrix eigenvalues,
Fourier comparisons for $n=5,6,7,8$ and collision offsets, exact rational
residual and isolating-interval checks, and deterministic finite-grid
checks. Tolerances and per-case data are recorded. Floating-point
corroboration is not interval certification, proof of openness, or an
independent audit; the module remains DRAFT.

\subsection{Priority checks for the independent P08 auditor}
\begin{enumerate}[label=(\arabic*),leftmargin=*]
\item Authenticate the P05 audited snapshot, P01--P05 body hashes,
A01--A05, and A05 E01 closure. Separate historical computational
provenance notes from the current audited proof inputs.
\item Check the full common domain, unique $x_+$ on $\mathcal C_+$,
and the explicitly fixed lower branch in the two-root construction.
Do not identify amplitudes by coefficient names.
\item Check positivity only on the maximal interval; separate sufficient
full-system bounds, scalar growth and blow-up, and necessary conditions.
Do not infer global boundedness merely from $M>0$.
\item Recompute the MB pre-equilibrium diagonal, all collision sums,
complex phases and imaginary signs, normality and both proofs of
$\lambda_0^{\MB}=-a$.
\item Check closed/strict selection cells, all maximal ties, parity,
finite-parameter endpoints, MB realizability, and each cell's exact
stability inequality. Generic P05 grid facts and DM realization are distinct.
\item Reconstruct the all-pair family, positivity, original residuals,
$\mu>0$, and both strict all-mode stability bounds.
\item Check the universal rational witness, finite-label bound
$p(u_J)\ge3/16$, and algebraic-integer exclusion of a tie at $2/9$.
\item Check the modal-gap obstruction and its branch restriction,
the all-$n\equiv3$ reverse family and $37/256$ estimate, and the
rational seventh-root isolating intervals and stability inequalities.
Independently check the $n=8$ both-unstable example.
\item Check the lower-root family and full-domain $4\mid n$ obstruction;
verify both distinct roots and that the compared branch was fixed in advance.
\item Verify physical five-dimensional openness using simple scalar
roots and finitely many strict spectral and representative gaps.
\item Confirm no use of P06, P07, P09--P18, no later bifurcation assertion,
no six-term substitution, and no claim that compilation verifies mathematics.
\end{enumerate}
